\section{Introducción}
\label{sec:introduccion}

En una fundación de hormigón masivo, la hidratación del cemento libera calor en todo el volumen, mientras que su disipación ocurre principalmente por las superficies. El núcleo puede seguir caliente cuando la cara exterior ya se está enfriando. En fundaciones de aerogeneradores, \citet{ju2019} observaron esta diferencia y utilizaron mediciones a distintas profundidades para evaluar el control térmico.

El plazo de colocación y el control térmico están relacionados. Colocar más despacio modifica las edades relativas de las capas y la hora a la que quedan expuestas al ambiente. A su vez, retirar la aislación cuando el núcleo sigue caliente puede ampliar el diferencial \citep{ichControlTermico}. En una fundación de puente, \citet{sargam2019} contrastaron las temperaturas calculadas con mediciones de obra y mostraron la sensibilidad de la evolución térmica a los materiales y a las condiciones constructivas.

Este estudio examina si es posible sostener la mayor tasa de colocación disponible sin exceder el riesgo térmico admitido. La hipótesis de trabajo es que una colocación continua puede cumplir los criterios docentes si la refrigeración y la protección de las superficies se mantienen durante las etapas críticas. Para evaluarla se compara cada plan bajo las mismas combinaciones de hidratación y clima.

\subsection{Objetivo general}
Identificar el plan de colocación de menor duración entre las alternativas consideradas para la fundación del caso docente en Los Ángeles, región del Biobío, manteniendo una probabilidad conjunta no mayor que el $5~\%$ de exceder una temperatura de $70\,^{\circ}\mathrm{C}$ o un diferencial térmico de $20\,\mathrm{K}$, bajo el modelo de incertidumbre adoptado. Estos límites corresponden a criterios académicos.

\subsection{Objetivos específicos}
\begin{itemize}
    \item Documentar la procedencia de los datos y los
    supuestos del modelo térmico.
    \item Verificar el método mediante casos conocidos,
    balances energéticos y análisis de sensibilidad
    espacial y temporal.
    \item Comparar la duración de colocación y las
    probabilidades de excedencia de los planes bajo
    escenarios comunes.
    \item Formular un plan de control térmico con
    monitoreo y acciones de contingencia.
\end{itemize}